\documentclass[10pt,a4paper]{amsart}
\usepackage[margin=19mm]{geometry}
\usepackage{amsmath,amssymb,mathtools}
\usepackage[scr=rsfs,cal=euler]{mathalfa}
\usepackage{xfrac}
\usepackage[unicode,hidelinks,bookmarks=false]{hyperref}
\newcommand{\ie}{i.e.\ }
\newcommand{\cf}{cf.\ }
\newcommand{\resp}{respectively\ }
\newcommand{\Subsection}{\subsection}
\providecommand{\nobreakdash}{\nobreak\hbox{-}}
\setlength{\emergencystretch}{2em}
\multlinegap0pt
\usepackage{enumerate}
\usepackage{amssymb}
\usepackage{tikz}
\usepackage{bm}
\usepackage{dsfont}
\newtheorem{proposition}{Proposition}
\def\block{{\rm blk}}
\title{Gradient Flow Equations for Deep Linear Neural Networks: A Survey from a Network Perspective}
\author{Joel Wendin, Claudio Altafini}
\date{Source: arXiv:2511.10362v1 (CC BY 4.0)}
\begin{document}
\maketitle

\noindent\emph{Source and licence.} Gradient Flow Equations for Deep Linear Neural Networks: A Survey from a Network Perspective. Version arXiv:2511.10362v1 (CC BY 4.0), distributed by arXiv under the Creative Commons Attribution 4.0 International licence, http://creativecommons.org/licenses/by/4.0/, as recorded in the arXiv licence metadata for this exact version and shown on the abstract page https://arxiv.org/abs/2511.10362v1. Full licence notice: \texttt{licenses/arxiv-2511.10362-CC-BY-4.0.txt}.\par\smallskip

\noindent\textit{Editorial scope.} Counted unit: the article's proposition prop:svdA (source0.tex lines 2052--2056). The article's complete original proof of it is delivered with it (source0.tex lines 2058--2151, one \texttt{proof} environment of 94 lines). No mathematics is written, repaired or re-derived: the statement and the proof are the article's own lines, in the article's own notation, and no retained line is substituted or re-flowed.  No further source-local context is needed: every cross-reference of every retained line resolves inside the two delivered spans.  Nothing was omitted from any retained span, so the package carries no omission record.  No retained line contains a citation, so the wrapper carries no bibliography and no bibliography entry.  No retained line contains a figure environment, an image-inclusion command or drawing code, so no image is delivered and the package declares no asset dependency.  Editorial formatting only, in addition: an AMS \texttt{amsart} preamble that restates the article's own declarations and macros used by the retained lines, namely the environments \texttt{proposition} on separate counters, together with the standard packages those lines need.  The retained environments print in this document as Proposition 1 in order of appearance, whereas the article prints the same items with the numbering of its own full document.  The complete native source of the article as distributed in the arXiv e-print for this exact version is bundled unchanged as \texttt{sources/189033/source0.tex}.
\begin{proposition}
\label{prop:svdA}
Let $ A= \block_1(W_1, \ldots , W_h)\in \mathcal{A} $. The singular values of $A$ are given by the union of the singular values of each block $W_i$, with $n_v-q$ extra zeros, where $q=\sum_{i=1}^{h}\min\{d_i,d_{i-1}\}$.
The associated singular vectors of $A$ can also be obtained assembling the singular vectors of the blocks $ W_i$.
\end{proposition}

\begin{proof}
Consider a real SVD of each block $W_i=U_i \Omega_iV_i^\top \in\mathbb{R}^{d_i\times d_{i-1}}$, $U_i\in\mathbb{R}^{d_i\times d_i}$, $\Omega_i\in\mathbb{R}^{d_i\times d_{i-i}}$, and $V_i\in\mathbb{R}^{d_{i-1}\times d_{i-1}}$. As $\Omega_i$ may be rectangular, we write $\Omega_i = \begin{bmatrix}
\Lambda_i & 0
\end{bmatrix}$ (or $\Omega_i = \begin{bmatrix}
\Lambda_i & 0
\end{bmatrix}^\top $) with $\Lambda_i$ diagonal. 
We divide the matrices into columns corresponding to singular values, and columns that have no contribution to $W_i$:
\begin{equation}
	U_i = \begin{bmatrix}
		U_i^\sigma & U_i^{o}
	\end{bmatrix}, \quad V_i = \begin{bmatrix}
		V_i^\sigma & V_i^{o}
	\end{bmatrix}
\end{equation}
with $U_i^\sigma \in\mathbb{R}^{d_i\times q_i}$, $U_i^{o}\in\mathbb{R}^{d_i\times p_i^u}$, $\Lambda_i\in\mathbb{R}^{q_i\times q_i}$, $V_i^\sigma \in\mathbb{R}^{d_{i-1}\times q_i}$, $V_i^{o}\in\mathbb{R}^{d_{i-1}\times p_i^v}$, where we have defined $q_i=\min\{d_i,d_{i-1}\}$, $p_i^u=d_i-q_i$ and $p_i^v=d_{i-1}-q_i$, accordingly, for $i=1,...,h$. Let also $q=\sum_{i=1}^h q_i$, $p^u=\sum_{i=1}^h p_i^u$, $p^v=\sum_{i=1}^h p_i^v$. Note that $W_i=U_i^\sigma \Lambda_iV_i^{\sigma \top}$, corresponding to a thin SVD. Now let
\begin{equation}
	U = 
	\left[
	\begin{array}{cccc|cccc|c}
		0 & & & & 0 & & & & U_{h+1} \\
		U_1^\sigma & 0 & & & U_{1}^o & 0 & & &0  \\
		0 & U_2^\sigma & \ddots & & 0 & U_{2}^o &\ddots  & & \vdots\\
		& \ddots & \ddots &0  & & \ddots & \ddots & 0 & \\
		& & 0 & U_h^\sigma & & & 0 & U_{h}^o & 0\\
	\end{array}
	\right] = \left[
	\begin{array}{c|c|c}
		U^\sigma & U^o & U_{h+1}^o \\
	\end{array}
	\right] \in\mathbb{R}^{n_v\times n_v}
    \label{eq:svd_U}
\end{equation}
where $U^\sigma\in\mathbb{R}^{n_v\times q}$, $U^o\in\mathbb{R}^{n_v\times p^u}$, and $U_{h+1}^o\in\mathbb{R}^{n_v\times d_x}$ is a tall matrix with a unitary block $U_{h+1}\in\mathbb{R}^{d_x\times d_x}$ at the top. Similarly,
\begin{equation}
	V = 
	\left[
	\begin{array}{cccc|cccc|c}
		V_1^{\sigma} & 0 & & & V_{1}^o & 0 & & & 0 \\
		0 & V_2^{\sigma} & \ddots & & 0 & V_{2}^o & \ddots & & \vdots\\
		& \ddots & \ddots & 0 & & \ddots & \ddots & 0 & \\
		& & & V_{h}^{\sigma} & & & & V_{h}^{o} & 0 \\
		& & & 0 & & & & 0 & V_{h+1}^{\sigma} \\
	\end{array}
	\right] = \left[
	\begin{array}{c|c|c}
		V^\sigma & V^o & V_{h+1}^o \\
	\end{array}
	\right] \in\mathbb{R}^{n_v\times n_v}
    \label{eq:svd_V}
\end{equation}
with  $V^\sigma\in\mathbb{R}^{n_v\times q}$, $V^o\in\mathbb{R}^{n_v\times p^v}$, $V_{h+1}^o\in\mathbb{R}^{n_v\times d_y}$, and $V_{h+1}\in\mathbb{R}^{d_y\times d_y}$. Define the diagonal matrix
\begin{equation}\label{eq:svddiag}
	\Omega = \begin{bmatrix}
		\Lambda_1 & & & & \\
		& \Lambda_2 & & & \\
		& & \ddots & & \\
		& & & \Lambda_h & \\
		& & & & 0 \\
	\end{bmatrix} = \begin{bmatrix}
	\Lambda & 0 \\
	0 & 0
	\end{bmatrix} \in\mathbb{R}^{n_v\times n_v}
\end{equation}
where the diagonal block $\Lambda$ is $q\times q$, and the last $n_v-q$ rows and columns of $\Omega$ are zeros. This means that $U\Omega V^\top  = U^\sigma \Lambda V^{\sigma \top}$, hence
\begin{align}
	U\Omega V^\top  & = \begin{bmatrix}
		0 & & & \\
		U_1^\sigma & 0 & & \\
		0 & U_2^\sigma & \ddots & \\
		& \ddots & \ddots & 0\\
		& & 0 & U_h^\sigma \\
	\end{bmatrix}
	\begin{bmatrix}
		\Lambda_1 & & & \\
		& \Lambda_2 & & \\
		& & \ddots & \\
		& & & \Lambda_h \\
	\end{bmatrix}
	\begin{bmatrix}
		V_1^{\sigma \top} & 0 & & &  \\
		0 & V_2^{\sigma \top} & \ddots & & \\
		& \ddots & \ddots & 0 & \\
		& & 0 & V_h^{\sigma \top} & 0 \\
	\end{bmatrix} \nonumber \\
	& =\begin{bmatrix}
		0 & & & & \\
		U_1^\sigma \Lambda_1 V_1^{\sigma \top} & 0 & & & \\
		0 & U_2^\sigma \Lambda_2 V_2^{\sigma \top} & \ddots & & \\
		& \ddots & \ddots & & \\
		& & 0 & U_h^\sigma \Lambda_h V_h^{\sigma \top} & 0 \\
	\end{bmatrix}
\end{align}
As noted before, $W_i=U_i^\sigma \Lambda_iV_i^{\sigma \top}$, so $A=U\Omega V^\top $. $\Omega$ is diagonal by construction, and each column in $U$ is orthogonal to all other columns in $U$ (this follows from each $U_i$ being orthogonal, and columns from different $U_i$:s in $U$ having no nonzero elements in the same position), this of course holds also for $V$. Thus, $A=U\Omega V^\top $ is a singular value decomposition of $A$, and it has the same singular values as the blocks, with $n_v-q$ extra zeros, as seen from \eqref{eq:svddiag}.
\end{proof}

\end{document}
